% from ega0-1-fr.tex
\oldpage[0\textsubscript{I}]{59}

  \item[$c)$] Le $A$-module $B$ est plat (autrement dit, le foncteur
    $M_{(B)}$ est exact), et pour tout $A$-module $M$, l'homomorphisme
    $x\mapsto1\otimes x$ de $M$ dans $M_{(B)}$ est injectif.
  \item[$d)$] Le $A$-module $B$ est plat et pour tout idéal $\mathfrak a$ de
    $A$, on a $\varphi^{-1}(\mathfrak aB)=\mathfrak a$.
  \item[$e)$] Le $A$-module $B$ est plat et pour tout idéal maximal
    $\mathfrak m$ de $A$, il existe un idéal maximal $\mathfrak n$ de $B$ tel
    que $\varphi^{-1}(\mathfrak n)=\mathfrak m$.
\end{enumerate}

Lorsque ces conditions ont lieu, on identifie $A$ à un sous-anneau de $B$.
\end{env}

\begin{env}[6.6.2]
\label{0.6.6.2-fr}
Soient $A$ un anneau local, $\mathfrak m$ son idéal maximal, $B$ une
$A$-algèbre telle que $\mathfrak mB\neq B$ (ce qui a lieu par exemple lorsque
$B$ est un anneau local et $A\to B$ un homomorphisme local). Si $B$ est un
$A$-module plat, $B$ est un $A$-module fidèlement plat. En effet, comme
$\mathfrak mB\neq B$, il y a un idéal maximal $\mathfrak n$ de $B$ contenant
$\mathfrak mB$~; comme $\varphi^{-1}(\mathfrak n)\cap A$ contient
$\mathfrak m$ et ne contient pas $1$, on a
$\varphi^{-1}(\mathfrak n)=\mathfrak m$ et le critère $e)$ de (6.6.1)
s'applique. Sous les conditions indiquées, on voit donc que si $B$ est
noethérien, il en est de même de $A$ (6.5.2).
\end{env}

\begin{env}[6.6.3]
\label{0.6.6.3-fr}
Soit $B$ une $A$-algèbre qui est un $A$-module fidèlement plat. Pour tout
$A$-module $M$ et tout sous-$A$-module $M'$ de $M$, on a (en identifiant $M$
à un sous-$A$-module de $M_{(B)}$) $M'=M\cap M'_{(B)}$. Pour que $M$ soit un
$A$-module plat (resp. fidèlement plat), il faut et il suffit que $M_{(B)}$
soit un $B$-module plat (resp. fidèlement plat).
\end{env}

\begin{env}[6.6.4]
\label{0.6.6.4-fr}
Soient $B$ une $A$-algèbre, $N$ un $B$-module fidèlement plat. Pour que $B$
soit un $A$-module plat (resp. fidèlement plat), il faut et il suffit que $N$
le soit.

En particulier, soit $C$ une $B$-algèbre~; si l'anneau $C$ est fidèlement
plat sur $B$ et $B$ fidèlement plat sur $A$, alors $C$ est fidèlement plat
sur $A$~; si $C$ est fidèlement plat sur $B$ et sur $A$, alors $B$ est
fidèlement plat sur $A$.
\end{env}

\subsection{Morphismes plats d'espaces annelés}
\label{subsection:0.6.7-fr}

\begin{env}[6.7.1]
\label{0.6.7.1-fr}
Soit $f:X\to Y$ un morphisme d'espaces annelés, et soit $\mathcal F$ un
$\mathcal O_X$-Module. On dit que $\mathcal F$ est $f$-plat (ou $Y$-plat
lorsque aucune confusion n'est à craindre sur $f$) en un point $x\in X$ si
$\mathcal F_x$ est un $\mathcal O_{f(x)}$-module plat~; on dit que
$\mathcal F$ est $f$-plat au-dessus de $y\in Y$ si $\mathcal F$ est $f$-plat
en tous les points $x\in f^{-1}(y)$~; on dit que $\mathcal F$ est $f$-plat
si $\mathcal F$ est $f$-plat en tous les points de $X$. On dit que le
morphisme $f$ est plat en $x\in X$ (resp. plat au-dessus de $y\in Y$, resp.
plat) si $\mathcal O_X$ est $f$-plat en $x$ (resp. $f$-plat au-dessus de $y$,
resp. $f$-plat).
\end{env}

\begin{env}[6.7.2]
\label{0.6.7.2-fr}
Avec les notations de (6.7.1), si $\mathcal F$ est $f$-plat en $x$, pour tout
voisinage ouvert $U$ de $y=f(x)$, le foncteur
$(f^*(\mathcal G)\otimes_{\mathcal O_X}\mathcal F)_x$ en $\mathcal G$ est
exact dans la catégorie des $(\mathcal O_Y|U)$-Modules~; en effet, cette
fibre s'identifie canoniquement à
$\mathcal G_y\otimes_{\mathcal O_y}\mathcal F_x$, et notre assertion résulte
de la définition. En particulier, si $f$ est un morphisme plat, le foncteur
$f^*$ est exact dans la catégorie des $\mathcal O_Y$-Modules.
\end{env}

\begin{env}[6.7.3]
\label{0.6.7.3-fr}
Inversement, supposons le faisceau d'anneaux $\mathcal O_Y$ cohérent, et
supposons que pour tout voisinage ouvert $U$ de $y$, le foncteur
$(f^*(\mathcal G)\otimes_{\mathcal O_X}\mathcal F)_x$ soit exact en
$\mathcal G$ dans la catégorie des $(\mathcal O_Y|U)$-Modules cohérents.
Alors $\mathcal F$ est $f$-plat en $x$. En effet, il suffit de prouver que
pour tout idéal de type fini $\mathfrak J$ de $\mathcal O_y$,
l'homomorphisme canonique
$\mathfrak J\otimes_{\mathcal O_y}\mathcal F_x\to\mathcal F_x$ est injectif
(6.1.1). Or, on sait (5.3.8) qu'il existe alors un voisinage

